% !TEX program = lualatex
\documentclass[12pt,a4paper]{ltjsarticle}
\usepackage{amsmath,amssymb,amsthm}
\usepackage{geometry}
\geometry{margin=2.5cm}
\usepackage{hyperref}
\hypersetup{colorlinks=true,linkcolor=blue,citecolor=blue,urlcolor=blue}
\usepackage{booktabs}
\usepackage{luatexja-fontspec}

% 定理環境
\newtheorem{theorem}{定理}[section]
\newtheorem{proposition}[theorem]{命題}
\newtheorem{lemma}[theorem]{補題}
\newtheorem{corollary}[theorem]{系}
\newtheorem{definition}[theorem]{定義}
\newtheorem*{remark}{注意}
\newtheorem*{observation}{観察}

\title{\textbf{離散グノモン正写像における電荷次元の\\幾何学的起源：5次元時空と電荷量子化}}

\author{木原 範昭（Noriaki Kihara）\\
\small WF System Co., Ltd.\\
\small （大阪大学 基礎工学部 卒業）\\
\small \href{https://orcid.org/0009-0004-6753-4020}{ORCID: 0009-0004-6753-4020}\\
\small \href{https://doi.org/10.5281/zenodo.19435162}{DOI: 10.5281/zenodo.19435162}}

\date{2026年4月\\[5pt]\small 研究ノート（幾何学的考察）}

\begin{document}

\maketitle

\begin{abstract}
前論文~\cite{Paper2} において，離散 Minkowski 格子 $\mathbb{Z}^{1+N}$ 上の等方的測地線方程式が $N=3$ では解を持たず（算術的アノマリー），$N=4$ で初めて解が存在することを示した．本稿は，この数論的制約が要請する5次元目を電荷（Coulomb）次元と同定し，グノモン正写像を $S^5(R)$ へ拡張する．

新たな結合定数は導入しない．唯一の結合定数 $C$（次元：長さ）を保持し，電荷次元と長さ次元を結ぶ変換係数 $Q = C/e_0$（$e_0$：最小電荷量子）を定義する．離散格子上の座標 $\tilde{\phi} = q/e_0 \in \mathbb{Z}$ は整数電荷数となり，電荷量子化はコンパクト化を仮定することなく離散格子の帰結として導出される．さらに，電荷方向のグノモン写像から $1/r^2$ Coulomb 法則を導出し，\cite{Paper2} の Lorentz 不変性と組み合わせることで Maxwell 方程式の全体を幾何学的に導出する．曲率方向の観測可能性の差異から，電磁波（スピン1）と重力波（スピン2）の偏向構造の違い，および電磁力に等価原理が成立しない理由が自然に導かれる．
\end{abstract}

\tableofcontents

%========================================
\section{はじめに：離散化前提の確認}
%========================================

\subsection{前提：離散グノモン写像は検証済み}

本稿は，前論文~\cite{Paper2} で確立された\textbf{離散グノモン正写像}を出発点とする．以下の結果は検証済みであり，本稿の前提として使用する：

\begin{enumerate}
\item 結合定数 $C$ により連続座標を離散化：$\tilde{x}^\mu = x^\mu/C \in \mathbb{Z}$．
\item 曲率パラメータ $w = C/R \in \mathbb{Z}_{\geq 0}$ のもとで離散グノモン写像 $\tilde{Y}^\mu = \tilde{x}^\mu/\tilde{L}$ を定義．
\item 全計量は2乗形式：$\tilde{\sigma}^2 = \sum_\mu \epsilon_\mu (\tilde{x}^\mu)^2$，$\tilde{L}^2 = 1 + w^2\tilde{\sigma}^2$．
\item Euclidean 版では $\tilde{L}^2 \geq 1 > 0$ が全格子点で保証（\cite{Paper2} 定理~6.1(i)）．
\end{enumerate}

\subsection{2乗計量の根本的理由：負の座標値の扱い}

離散格子 $\mathbb{Z}^n$ では各座標軸が正負の整数値をとる．計量は符号不変でなければならない：$(-n)^2 = n^2$．2乗形式はこれを保証する最小構造である：
\begin{equation}
\text{離散格子} \;\Longrightarrow\; \text{負の座標値} \;\Longrightarrow\; \text{符号不変性が必要} \;\Longrightarrow\; (\tilde{x}^\mu)^2 \tag{1.1}
\end{equation}

\subsection{数論的動機付け：なぜ5次元が必要か}

\cite{Paper2} 注意~6.2 および~\cite{Kihara2026b} より：$\mathbb{Z}^{1+N}$ 上の等方測地線方程式 $l^2 + t^2 = Nx^2$ が解を持つ必要十分条件は，$N$ の素因数分解において $p \equiv 3 \pmod{4}$ の素数が偶数乗であること（Fermat）．$N=3$ はこれに違反；$N=4$ は充足．\textbf{離散等方性には $N=4$ 以上の空間次元が必要．}

\subsection{電荷次元の観測可能性}

\cite{Paper2} 公理~3.1 では，$R$ 方向の不可観測性が虚軸化（Lorentzian 符号）を導いた．電荷次元 $\phi$ は\textbf{観測可能}（電荷は測定できる）であるため，公理~B は\textbf{適用されない}：$\epsilon_4 = +1$．

\begin{center}
\begin{tabular}{lll}
\toprule
次元 & 観測可能性 & 公理~B（虚軸化）の適用 \\
\midrule
$Y^{n+1}$（$R$ 方向） & 不可観測 & 適用 → 符号反転 \\
$\phi = Qq$（電荷方向） & \textbf{観測可能} & \textbf{非適用} → 符号不変 \\
\bottomrule
\end{tabular}
\end{center}

\subsection{本稿の構成}

\S2：5次元グノモン写像．\S3：変換係数 $Q$ と電荷量子化．\S4：5次元離散写像．\S5：離散等方性の回復．\S6：5次元 Einstein テンソル．\S7：整合性確認．\S8：今後の課題と電磁場．

%========================================
\section{5次元グノモン写像の拡張}
%========================================

\subsection{座標系の定義}

\begin{equation}
(x^1, x^2, x^3, x^4, x^5) = (x, y, z, \phi, ct) \tag{2.1}
\end{equation}
符号：$\epsilon_\mu = (+1,+1,+1,+1,-1)$．

\subsection{5次元グノモン写像}

\cite{Paper1} を $n=5$ に適用：
\begin{equation}
\Phi: \Pi_R \to S^5(R), \quad Y^\mu = \frac{Rx^\mu}{l}, \quad Y^6 = \frac{R^2}{l} \tag{2.3}
\end{equation}
\begin{equation}
l^2 = R^2 + x^2 + y^2 + z^2 + Q^2q^2 - c^2t^2 \tag{2.4}
\end{equation}
引き戻し計量：
\begin{equation}
g_{\mu\nu} = \frac{R^2}{l^2}\left(\epsilon_\mu \delta_{\mu\nu} - \frac{\epsilon_\mu\epsilon_\nu x_\mu x_\nu}{l^2}\right) \tag{2.5}
\end{equation}

%========================================
\section{変換係数 $Q$ と電荷量子化}
%========================================

\subsection{$c$ との構造的対称性}

光速 $c$ は時間を長さに変換し，$Q$ は電荷を長さに変換する：

\begin{center}
\begin{tabular}{lll}
\toprule
 & 時間 → 長さ & 電荷 → 長さ \\
\midrule
変換係数 & $c$ & $Q$ \\
変換式 & $x^5 = ct$ & $x^4 = Qq$ \\
次元 & m/s & m/C \\
\bottomrule
\end{tabular}
\end{center}

\subsection{$Q$ の決定}

離散化条件 $\tilde{\phi} = Qq/C \in \mathbb{Z}$ は最小電荷量子 $e_0$ の存在を要請する：
\begin{equation}
Q = \frac{C}{e_0}, \qquad \tilde{\phi} = \frac{q}{e_0} \in \mathbb{Z} \tag{3.3}
\end{equation}
$e_0 = e/3$（クォーク電荷）として：
\begin{equation}
Q = \frac{3C}{e} \approx 3.03 \times 10^{-16}\;\text{m/C} \tag{3.5}
\end{equation}

\subsection{電荷量子化}

\begin{theorem}[電荷量子化]\label{thm:charge}
$\tilde{\phi} \in \mathbb{Z}$ および $Q = C/e_0$ のもとで：
\begin{equation}
q = n \cdot e_0, \qquad n \in \mathbb{Z} \tag{3.7}
\end{equation}
\end{theorem}

\begin{remark}[コンパクト化は不要]
Kaluza--Klein 理論~\cite{Kaluza1921,Klein1926} と異なり，電荷量子化は離散格子 $\tilde{\phi} \in \mathbb{Z}$ から直接導出される．$\phi$ 方向は非コンパクト．
\end{remark}

\begin{remark}[既知粒子の離散座標]
$e_0 = e/3$ として：陽子（$\tilde{\phi}=+3$），電子（$-3$），u クォーク（$+2$），d クォーク（$-1$），ニュートリノ（$0$）．すべて整数．
\end{remark}

%========================================
\section{5次元離散グノモン写像}
%========================================

\subsection{離散化}

\begin{equation}
\tilde{x} = x/C,\; \tilde{y} = y/C,\; \tilde{z} = z/C,\; \tilde{\phi} = q/e_0,\; \tilde{t} = ct/C \;\in \mathbb{Z} \tag{4.1}
\end{equation}

\subsection{離散計量}

\begin{equation}
\tilde{L}^2 = 1 + w^2\tilde{\sigma}^2, \qquad \tilde{\sigma}^2 = \tilde{x}^2 + \tilde{y}^2 + \tilde{z}^2 + \tilde{\phi}^2 - \tilde{t}^2 \tag{4.2}
\end{equation}

\subsection{正則性（5次元版）}

\begin{theorem}[5次元正則性]\label{thm:5dreg}
\textbf{(i)} Euclidean 符号：$\tilde{L}^2 \geq 1 > 0$（大域的正則）．

\textbf{(ii)} Lorentzian 符号：ホライズンは $c^2t^2 = R^2 + x^2 + y^2 + z^2 + Q^2q^2$．電荷項 $Q^2q^2 \geq 0$ により正則領域が4次元版より\textbf{拡大}．
\end{theorem}

%========================================
\section{離散等方性の回復（$N=4$）}
%========================================

$N=4$：$4 = 2^2$ には $p \equiv 3 \pmod{4}$ の素因数なし．方程式 $l^2 + t^2 = 4x^2$ は解を持つ．最小解：$l=0, t=2, x=1$．

\begin{proposition}[離散等方性の回復]\label{prop:isotropy}
$\mathbb{Z}^{1+3}$ には等方的 null 測地線が存在しないが，電荷次元を加えた $\mathbb{Z}^{1+4}$ には存在する．電荷次元は離散等方性のために数論的に必要．
\end{proposition}

%========================================
\section{5次元 Einstein テンソル}
%========================================

\cite{Paper1} 命題~2.1 を $n=5$ に適用：
\begin{equation}
\Lambda_5 = \frac{6}{R^2}, \qquad G_{\mu\nu} + \frac{6}{R^2}g_{\mu\nu} = 0 \tag{6.1--6.2}
\end{equation}

%========================================
\section{整合性の確認}
%========================================

\textbf{結合定数：}$C$ のみ．$Q = C/e_0$ は変換係数であり自由パラメータではない．

\begin{center}
\begin{tabular}{llll}
\toprule
物理量 & 変換 & 離散座標 & 量子化 \\
\midrule
空間位置 & $x$ [m] & $\tilde{x} = x/C \in \mathbb{Z}$ & 長さ \\
時間 & $t$ [s] → $ct$ [m] & $\tilde{t} = ct/C \in \mathbb{Z}$ & 時間 \\
電荷 & $q$ [C] → $Qq$ [m] & $\tilde{\phi} = q/e_0 \in \mathbb{Z}$ & \textbf{電荷} \\
\bottomrule
\end{tabular}
\end{center}

\begin{proposition}[負座標値の安全性]\label{prop:safety}
離散グノモン写像において，$\tilde{L}^2$ は全 $\tilde{x}^\mu$ の符号変化に対して不変．
\end{proposition}

%========================================
\section{今後の課題と電磁場}
%========================================

\subsection{微細構造定数 $\alpha$ の導出}

$\alpha \approx 1/137$ の理論的導出は本稿では与えない．Wyler 公式~\cite{Wyler1971} と $S^5(R)$ の対称群 $SO(5,1)$ の接続が重要課題．

\subsection{電荷方向の電磁場}

\subsubsection{$\phi$ 方向の有向曲率}

\begin{center}
\begin{tabular}{lll}
\toprule
 & 重力 & 電磁力 \\
\midrule
曲率方向 & $R$（不可観測） & $\phi = Qq$（\textbf{観測可能・有向}） \\
曲率半径 & $R$ & $\phi_0 = Qq_0$（ソース電荷） \\
性質 & 一様 & \textbf{有向}（$q_0$ の符号） \\
\bottomrule
\end{tabular}
\end{center}

\subsubsection{$1/r^2$ 法則の導出}

ソース電荷 $q_0$ が $(x,y,z,t)$ 空間を曲率半径 $\phi_0 = Qq_0$ で曲げる．グノモン写像の埋め込み座標 $Y^i = \phi_0 x^i / l_q$ の時間二階微分（$t=0$）：
\begin{equation}
\frac{d^2 Y^i}{dt^2}\bigg|_{t=0} = \frac{\phi_0\, r_i\, c^2}{(\phi_0^2 + r^2)^{3/2}} \tag{8.7}
\end{equation}
$r \gg |\phi_0|$ の極限：
\begin{equation}
\boxed{\frac{d^2 Y^i}{dt^2} \approx \frac{Qq_0 \cdot c^2}{r^2}\,\hat{r}_i} \tag{8.8}
\end{equation}
\textbf{これは $1/r^2$ Coulomb 法則である．}

\subsubsection{テスト粒子応答と等価原理の不在}

\begin{remark}[等価原理の不在の幾何学的説明]
重力では $R$ が不可観測であるため全粒子が同一に応答（等価原理）．電磁力では $\phi$ が観測可能であるため応答は $q_1$ に依存（等価原理の不在）．等価原理は $R$ の不可観測性の\textbf{帰結}である．
\end{remark}

\subsubsection{偏向構造：スピン1対スピン2}

\begin{proposition}[偏向と観測可能性]\label{prop:polar}
電磁波がスピン1（ベクトル偏向）であり，重力波がスピン2（テンソル偏向）であることは，それぞれの曲率方向の観測可能性の差異に帰着する．
\end{proposition}

\subsubsection{Maxwell 方程式の完全形}

\begin{proposition}[Maxwell 方程式の幾何学的導出]\label{prop:maxwell}
本稿の Coulomb 法則~(8.8) と \cite{Paper2} の Lorentz 不変性を組み合わせることで，Maxwell 方程式の全体が幾何学的に導出される．
\end{proposition}

\begin{center}
\begin{tabular}{ll}
\toprule
Maxwell 方程式 & 導出 \\
\midrule
$\nabla \cdot \mathbf{E} = \rho/\epsilon_0$ & 式~(8.8) \\
$\nabla \cdot \mathbf{B} = 0$ & $\phi$ は1次元 → 磁気単極子の自由度なし \\
$\nabla \times \mathbf{E} = -\partial\mathbf{B}/\partial t$ & Lorentz 不変性~\cite{Paper2} \\
$\nabla \times \mathbf{B} = \mu_0\mathbf{J} + \mu_0\epsilon_0\partial\mathbf{E}/\partial t$ & Lorentz 不変性~\cite{Paper2} \\
\bottomrule
\end{tabular}
\end{center}

%========================================
\section*{結果の整理}
%========================================

\begin{center}
\begin{tabular}{llll}
\toprule
項目 & 結果 & 式番号 & 状態 \\
\midrule
数論的動機付け & $N\!=\!3$ 解なし；$N\!=\!4$ 解あり & \S1.3 & Fermat~\cite{Kihara2026b} \\
電荷の観測可能性 & $\phi$ 観測可能 → $\epsilon_4\!=\!+1$ & \S1.4 & 公理 A,B \\
変換係数 $Q$ & $Q = C/e_0$ & (3.3) & 定義 \\
電荷量子化 & $q = ne_0$ & (3.7) & 定理~\ref{thm:charge} \\
5次元正則性 & $\tilde{L}^2 \geq 1$；ホライズン拡大 & (4.4) & 定理~\ref{thm:5dreg} \\
離散等方性 & $t = 2x$（$N\!=\!4$） & (5.2) & 命題~\ref{prop:isotropy} \\
Coulomb 法則 & $d^2Y^i/dt^2 \approx Qq_0 c^2/r^2$ & (8.8) & \S8.2 \\
等価原理 & $R$ 不可観測 → あり；$\phi$ 観測可能 → なし & 注意~8.3 & \S8.2 \\
偏向構造 & EM: スピン1；GW: スピン2 & 命題~\ref{prop:polar} & \S8.2 \\
Maxwell 方程式 & Coulomb + Lorentz → 完全導出 & 命題~\ref{prop:maxwell} & \S8.2 \\
\bottomrule
\end{tabular}
\end{center}

\begin{thebibliography}{99}

\bibitem{Paper1}
Kihara, N. (2026). Geometric Formulation of 4-Dimensional Spacetime via Gnomonic Projection.
DOI: \href{https://doi.org/10.5281/zenodo.19427780}{10.5281/zenodo.19427780}.

\bibitem{Paper2}
Kihara, N. (2026). Geometric Origin of Lorentzian Spacetime via Gnomonic Projection: A Unified Description through the Discretization Constant $C$.
DOI: \href{https://doi.org/10.5281/zenodo.19434932}{10.5281/zenodo.19434932}.

\bibitem{Kihara2026b}
Kihara, N. (2026). 離散ミンコフスキー間隔におけるガウス整数を用いた正整数解の解法. GitHub: ai-chat-logs-open.

\bibitem{Kaluza1921}
T.~Kaluza, ``Zum Unit\"{a}tsproblem der Physik,'' \textit{Sitzungsber.\ Preuss.\ Akad.\ Wiss.}, 966--972, 1921.

\bibitem{Klein1926}
O.~Klein, ``Quantentheorie und f\"{u}nfdimensionale Relativit\"{a}tstheorie,'' \textit{Z.\ Phys.}, \textbf{37}, 895, 1926.

\bibitem{Wyler1971}
A.~Wyler, ``L'espace sym\'{e}trique du groupe des \'{e}quations de Maxwell,'' \textit{C.\ R.\ Acad.\ Sci.\ Paris}, \textbf{272}, 186, 1971.

\bibitem{Hurwitz1898}
A.~Hurwitz, ``\"{U}ber die Composition der quadratischen Formen von beliebig vielen Variablen,'' \textit{Nachr.\ Ges.\ Wiss.\ G\"{o}ttingen}, 309--316, 1898.

\end{thebibliography}

\bigskip
\noindent\textit{本稿の議論は純粋幾何学と数論の範囲に限定される．電磁場の力学，微細構造定数の導出，および既存理論との同定は今後の研究課題である．}\\
\textit{最終更新：2026-04-06}

\end{document}
